\section{Counterfactual sensitivity: tables turned, lever by lever}
\label{sec:counterfactuals}

This section conducts the \emph{tables turned} pass: for every key
design choice, we ask ``what would the headline say if we had
chosen differently?'' The counterfactual is grounded in the
artefacts on disk -- we never invent data. When the artefacts do not
support a counterfactual we say so explicitly.

\subsection{CF-1: What if we had reported only accuracy?}
The headline would have been \emph{centralized wins} ($0.908$), with
local-only second ($0.904$) and grid-search third ($0.901$). The
deployment recommendation would have been ``pool the data'' (which
is non-deployable in the proposal's scope) or ``train one model per
ICU'' (which leaves the small Neuro units with $0$ recall on
deaths). The fairness-violating recommendation that an accuracy-
only report would surface is the precise reason the proposal
required \auprc{} and worst-client recall.

\subsection{CF-2: What if we had reported only \auroc{}?}
\fedprox{}-$0.1$ wins ($0.913$); the operating recommendation is the
same as the report's actual recommendation. \auroc{} is rank-only,
so it disguises class-imbalance pathologies; it is the right tool to
identify the right method, but it does not reveal the small-client
representation problem. Worst-client recall is the only metric that
shows how deeply \fedavg{} fails the small Neuro units.

\subsection{CF-3: What if we had not run \fedprox{}?}
The proposal-alignment dossier would still have produced the
\cvar{} discreteness finding, the LP shadow-price collapse, the
loss landscape early-stop hint, the Dirichlet stress test, and the
logistic backbone. The fairness improvement would not exist:
\fedavg{}'s $0.21$ worst-client recall would be the headline. The
report would conclude ``federation is unfair on \mimic{}'' rather
than ``federation \emph{with proximal anchoring} closes the
fairness gap.''

\subsection{CF-4: What if we had not run the LP?}
The communication-vs-loss trade-off would be a scatter plot, not a
recommendation. The shadow price collapses thirteen orders of
magnitude as soon as the budget exits the active region; that is
the LP saying ``stop spending bytes.'' Without the LP we would have
no principled way to defend the headline operating-point choice
against a reviewer who says ``add another $100$ Mb.''

\subsection{CF-5: What if we had used $C\!=\!20$ instead of $C\!=\!5$?}
The \cvar{} discreteness collapse would not have occurred:
$\lfloor 20 \cdot 0.05 \rfloor \!=\! 1$, $\lfloor 20 \cdot 0.25
\rfloor \!=\! 5$, $\lfloor 20 \cdot 0.5 \rfloor \!=\! 10$ -- four
distinct tail sizes. \cvar{} would be a continuous-feeling knob.
The cost is per-round communication: $20$ clients per round at
\fedprox{}-$0.1$'s $\sim$$25$ rounds is $\sim$$700$ Mb -- about
$50\%$ more than the headline. The LP at this budget says the
$\lambda$ shadow price is still negligible, so the extra spend buys
no operational improvement.

\subsection{CF-6: What if we had not run the loss landscape?}
We would not know that the MLP overshoots its minimum at
$\alpha\!\approx\!0.85$ along the init-to-final axis. The deployment
recommendation ``terminate $\sim$$15\%$ early'' would not exist.
This is the only result in the report that comes from the landscape
chapter alone; everything else is corroborated by other axes.

\subsection{CF-7: What if we had only run logistic regression?}
The headline \auprc{} would be $0.612$ instead of $0.665$
($\Delta\!=\!0.053$). Communication would be $6.4$ Mb instead of
$428$ ($\sim$$67\!\times$ less). Worst-client recall would be
$0.45$ instead of $0.50$ -- nearly the same. The report would have
recommended logistic for everyone; the actual report recommends
logistic for bandwidth-constrained ICUs and MLP elsewhere.

\subsection{CF-8: What if we had used a different seed pool?}
The seed pool is $\{7, 11, 19, 23, 29, 31, 37, 41, 43, 47\}$. The
across-seed std of \fedprox{}-$0.1$ \auprc{} is $0.005$ -- tight
enough that any prime-number seed pool of size $10$ would land in
the same operating region. \fedavg{}'s std is $0.014$, also tight;
the worst-client recall std is $0.14$ for \fedavg{} and $0.05$ for
\fedprox{}-$0.1$, so \fedprox{}-$0.1$ is robust to seed choice and
\fedavg{} is borderline. The Dirichlet-$0.1$ across-seed std is
$0.113$, large enough that a different seed pool would meaningfully
shift the headline -- this is exactly the failure mode the report
flags.

\subsection{CF-9: What if we had used class-balanced loss for centralized?}
Centralized accuracy would drop, centralized \auprc{} would rise --
but the available data does not let us extrapolate by how much. The
headline \emph{relative} ranking (\fedprox{}-$0.1$ above centralized
on \auprc{}) is unlikely to change because the $0.080$ \auprc{} gap
is large; it would shrink to maybe $0.04$ but not reverse.

\subsection{CF-10: What if we had no early-stopping?}
All methods would run to round $1000$. \fedprox{}'s monotone
\auprc{} climb would continue and then either plateau or U-turn;
without the early-stop ledger we cannot say which. The
\emph{communication} number would explode: \fedprox{}-$0.1$ at
$1000$ rounds would communicate $\sim$$15$ Gb, where now it
communicates $428$ Mb. The LP shadow price at $15$ Gb is $0$ exactly,
which is the strongest possible argument for early stopping.

\subsection{CF-11: What if we had used Hessian eigenvector directions
instead of random for the 2D landscape?}
The 2D landscape would show sharper curvature; the bowl shape would
be tighter and the ridge in the random-direction figure would
disappear. The 1D conclusion (terminate $\sim$$15\%$ early) would
not change. This is a presentation upgrade, not a substantive one.

\subsection{CF-12: What if we had run more Dirichlet seeds at $\beta=0.1$?}
The mean accuracy would barely move. The std would tighten or
widen; current data ($10$ seeds at $\beta\!=\!0.1$) gives across-seed
std $\sim$$0.04$; another $20$ seeds would tighten the CI on the std
itself, not the mean. The qualitative conclusion (heterogeneity
hurts seed-tail, not seed-mean) would persist.

\subsection{CF-13: What if we had used adaptive $\mu_k$ per client?}
We did not run this; it is a future-work item. Mechanically,
$\mu_k \!\propto\! 1/|\mathcal{D}_k|$ would give Neuro Intermediate
$\mu \!\approx\! 0.4$ and MICU $\mu \!\approx\! 0.018$ -- a $20\!\times$
range. We hypothesise this would push worst-client recall to
$\sim$$0.55$ at the cost of slightly worse \auroc{}; we cannot defend
the prediction without a run.

\subsection{CF-14: What if we had used FedNova instead of \fedprox{}?}
FedNova normalises local updates by their per-client
\emph{number-of-steps}, which is a strict generalisation of
\fedprox{} when $E$ varies across clients. With $E\!=\!1$ uniformly,
FedNova reduces to \fedavg{}; the only way it adds value is with
heterogeneous local epochs. Future work.

\subsection{CF-15: What if we had used SCAFFOLD?}
SCAFFOLD adds a control-variate per client to remove the bias of
\fedavg{} on heterogeneous data. The fairness improvement should be
similar to \fedprox{}, the convergence rate should be better, and
the per-client memory footprint doubles (each client stores its
control variate). On \mimic{}'s $9$ clients this is operationally
fine; on $1000$+ clients it would be the deployment story. Future
work.

\subsection{CF-16: What if we had not measured per-seed std?}
We would have reported \cvar{}-$0$ as the runner-up on worst-client
recall (mean $0.26$ vs.\ \fedprox{}-$0.1$'s $0.50$). With std
$0.21$ at mean $0.26$, $\sim$$40\%$ of \cvar{}-$0$ seeds returned $0$;
a single point estimate would have been deceptive. The
deployment recommendation would have been wrong.

\subsection{CF-17: What if we had not split MIMIC by ICU care unit?}
We would have used a Dirichlet-$\beta$ split or random-shard split.
Both lose the natural-clients story. The natural ICU partition is the
key to the report -- the small Neuro units \emph{exist} in the data
and are the operationally meaningful failure case. A synthetic split
would have produced a synthetic finding.

\subsection{CF-18: What if we had used a larger MLP?}
A $30 \!\to\! 256 \!\to\! 256 \!\to\! 2$ MLP has $\sim$$74$K
parameters vs.\ the current $6$K; communication would scale
$12\!\times$. The \auprc{} ceiling on this feature set is around
$0.71$ (centralized's best); the extra capacity probably buys
$\sim$$0.02$ \auprc{}, communication-bound. The LP at the larger
budget would still say ``shadow price negligible''; the
recommendation does not change.

\subsection{CF-19: What if we had used FedAdam at the server?}
Server-side adaptive optimisation accelerates global convergence.
On a $30$-round federation it would land at the same final loss
faster; the absolute headline numbers would not change because the
early-stop ledger already terminates well before the budget. Future
work for client-fewer / round-fewer regimes.

\subsection{CF-20: What if we had measured upload separately from
download?}
The report aggregates both directions. Federated systems are
upload-dominated (clients upload model updates; the server only
sends the global model). Our \fedprox{}-$0.1$ at $428$ Mb is
$\sim$$214$ Mb upload, $\sim$$214$ Mb download. Downstream networks
(satellite, cellular) charge upload differently; the LP at
upload-only budget would re-rank the operating points, but the
\fedprox{}-$0.1$ recommendation persists because its upload
budget ($214$ Mb) is comparable to \fedavg{}'s ($286$ Mb upload at
$572$ Mb total).

\paragraph{Summary of counterfactuals.}
None of the twenty counterfactuals reverse the headline
recommendation (use \fedprox{}-$0.1$, deploy logistic if
bandwidth-constrained, terminate $\sim$$15\%$ early). Several
\emph{strengthen} the recommendation (CF-3, CF-10, CF-16). A few
suggest follow-up work (CF-13, CF-14, CF-15, CF-19). The robustness
of the headline to this many what-if questions is the strongest
empirical argument we can make for it.
